\section{Evaluation}
\label{sec:evaluation}

We evaluate whether the mechanisms hold as described, in three ways:
automated tests that exercise each promise the way a person, an agent,
or an attacker would (\cref{sec:eval-suites,sec:eval-hostile});
a structured analysis of what each adversary can reach
(\cref{sec:eval-analysis}); and a check that the documentation this
paper is built on matches the code (\cref{sec:eval-docs}). We close
with what is not evaluated (\cref{sec:eval-not}).

\subsection{Test suites}
\label{sec:eval-suites}

\Cref{tab:suites} summarizes the suites. Runtime tests run real agent
code in real processes, with a \emph{scripted model}: a stand-in that
returns replies written in advance. This is deliberate. The claims
under test concern what the \emph{platform} does with whatever the
model proposes (G1), and a test must give the same result on every run.
The runtime suite runs against in-memory services that answer the
backend's contract; the cross-component suite includes an end-to-end
test in which a person's socket reaches the backend, the backend dials
a real runtime process, and the runtime calls back through the gateway
under the chat's delegation, with nothing stubbed but the model.

\begin{table}[t]
\centering\small
\caption{Automated test suites (test functions; parametrization makes
the executed count higher). Results from runs on 15~October~2026.}
\label{tab:suites}
\begin{tabularx}{\columnwidth}{@{}>{\raggedright\arraybackslash}p{0.28\columnwidth}r>{\raggedright\arraybackslash}X@{}}
\toprule
\textbf{Suite} & \textbf{Tests} & \textbf{Exercises} \\
\midrule
Runtime & 1{,}031 & assistant loop, executor gates, approvals, clock, stops, proxy, evidence, documentation \\
Backend & 924 & identity, policy invariants, delegation fence, data layer and encryption, OAuth, webhooks, audit \\
Cross-component & 180 & approval-to-package path; end-to-end chat across backend and runtime \\
Example agent & 24 & an agent run as its author would \\
\midrule
\textbf{Total} & \textbf{2{,}159} & \\
\quad of which sandbox, live & 83 & hostile agent in the deployed two-container topology (\cref{sec:eval-hostile}) \\
\bottomrule
\end{tabularx}
\end{table}

\paragraph{Results}
On the main branch on 15~October~2026 (Python 3.11, MongoDB~7), every
suite passed with no failures: runtime, 1{,}056 passed and 85 skipped
(the 83 live sandbox tests, which need the Linux image and are reported
in \cref{sec:eval-hostile}, and two others); backend, 969 passed and 6
skipped; cross-component, 216 passed and 3 skipped; example agent, 24
passed. Executed counts exceed \cref{tab:suites} because of
parametrization.

\paragraph{Regression tests from review}
In October 2026 the entire code base was reviewed line by line. Every
case found to make the code fail, hang, or answer wrongly was fixed
together with a test that repeats the attempt; in the runtime these
are kept in dedicated ``edge'' files---around stops (26 tests), the
executor (17), contracts and the SDK (31), the clock (8), session
restarts (6), concurrent stores (8), reasoning budgets (14), and
services and connectors (24). Design questions the review raised but
did not settle are published with the code (\cref{sec:limitations}).

\subsection{Confinement against a hostile agent}
\label{sec:eval-hostile}

The way to test a lock is to try to open it. The live sandbox suites
(\file{ai\_runtime/tests/test\_confinement\_live.py},
\file{test\_spawner\_live.py}) install a fixture agent that is
\emph{hostile on purpose}: each of its functions attempts one escape,
and the test passes only when the attempt is refused. Beside each
refusal, positive controls check that legitimate work still functions,
since a sandbox that refuses everything would pass the first kind of
test and fail the second. A second hostile package attacks at
\emph{build} time, from a local package index.

\paragraph{Setup}
We ran both files in the deployed topology: the runtime's image as the
runtime, and again as the agents' container beside it on an internal
network, with \code{NET\_ADMIN} for the egress rule. The host kernel was
Linux 6.18 (WSL2), which offers every Landlock feature in
\cref{tab:landlock}, so no test was skipped for want of a kernel
feature.

\paragraph{Results}
All 83 tests passed and none was skipped (58 confinement tests and 25
tests of the spawner across the two containers; 88\,s wall time).
\Cref{tab:hostile} groups the attempts by the promise they test and
the layer that refuses them.

\begin{table*}[t]
\centering\small
\caption{Escape attempts by the hostile agent, each refused, and the
legitimate operations verified alongside. All 83 tests passed on Linux
6.18 with no skips.}
\label{tab:hostile}
\begin{tabularx}{\textwidth}{@{}>{\raggedright\arraybackslash}p{0.15\textwidth}>{\raggedright\arraybackslash}X>{\raggedright\arraybackslash}p{0.2\textwidth}>{\raggedright\arraybackslash}p{0.24\textwidth}@{}}
\toprule
\textbf{Promise} & \textbf{Attempts (all refused)} & \textbf{Refused by} & \textbf{Still works} \\
\midrule
P1 own files & read the runtime's settings; write into its own package, environment, or the install directory; read or write another agent's home or package; read the uid table or the platform's code; write where every user may write & uid ownership; Landlock & reads what it runs from; spool shared with the runtime; home emptied each start, even if closed or deeply nested \\
P2 own user & signal the runtime; use the spawn helper to become another agent's uid; close or deeply nest its home to keep state across restarts & uid separation; \code{no\_new\_privs}; empty capability set; the helper's \code{clear} & runs as its own uid in its own home; stopping it ends everything it started \\
P3 declared hosts & connect past the proxy; reach the runtime's own door or the spawner; resolve a name directly; reach an undeclared host; reach anything with \code{hosts: []} or no network block; reach a blocked site with \code{hosts: any} & nftables owner rule; proxy; Landlock TCP rule & reaches its declared host; a host its credential names, only for that call; a host a person allowed on a code card, only until the call ends \\
P4 bounded resources & (verified, not attacked) the worker finds on itself exactly the process, open-file and file-size limits the container set, never a manifest's & \code{setrlimit} & threads and subprocesses start \\
P5 system calls & open an unused socket family; \code{ptrace}; create a namespace; connect to another agent's abstract Unix socket & seccomp; Landlock scoping (ABI~6) & its own Unix socket; a program it runs is fenced with it \\
Builds & a hostile build reads settings, writes into the environment, reaches hosts other than the index, fetches packages from elsewhere, or never ends & builder uid; proxy pass; five-minute kill & wheels built and installed by the runtime; only wheels are left \\
Records & an agent reads the monitoring record & folder mode and Landlock & both containers write one record \\
\bottomrule
\end{tabularx}
\end{table*}

\subsection{Security analysis}
\label{sec:eval-analysis}

\Cref{tab:matrix} relates each adversary to what it can do and what
holds it. Two observations follow. First, the defences against A1 and
A2 are the same set of gates, and none of them reads model output: this
is what makes the design independent of whether injection is detected.
Second, the defences sit in components of different trust: compromising
a less trusted component (a worker, then the runtime) leaves in place
every defence held by a more trusted one---the backend's keys, its
database, its decision on approvals, and its ability to withdraw
delegations.

\begin{table*}[t]
\centering\small
\caption{Adversaries, their capabilities, and the mechanisms that hold them.}
\label{tab:matrix}
\begin{tabularx}{\textwidth}{@{}>{\raggedright\arraybackslash}p{0.15\textwidth}>{\raggedright\arraybackslash}p{0.27\textwidth}>{\raggedright\arraybackslash}X@{}}
\toprule
\textbf{Adversary} & \textbf{Can} & \textbf{Held by} \\
\midrule
A1 model & emit any action in the vocabulary, with any inputs; write any prose & every action validated and gated outside the model; grants deny by default; the level gate; evidence drawn only from the trace \\
A2 injector & choose text the model reads, and so steer its proposals & the same gates; none of grants, level, manifest hosts, approval hash or credential binding derives from model output; a denial refuses the named hosts to the whole agent for the rest of the request \\
A3 agent author & arbitrary code in the worker during a call & approval pinned to commit and digest; per-operation resource mediation; uid, Landlock, seccomp, limits; egress only to declared hosts; dependencies built by a separate confined user \\
A4 another person & their own session, cookie, or API key & sessions resolved per request; deny-by-default policy; reads filtered by visibility inside the query; tenant scoping; \code{Origin} check on cookie-bearing sockets \\
A5 compromised worker & the above, plus attacks on co-tenants of the container & sandbox layers; internal network leading only to the runtime \\
A6 compromised runtime & every delegation it currently holds & no database route, decryption key, or backend signing key; 40-action fence; cannot decide approvals; stop-everything withdraws its delegations \\
\bottomrule
\end{tabularx}
\end{table*}

\paragraph{Injection reach by trust level}
Because the gate is a single comparison, the reach of a successful
injection (A2) inside one chat can be stated exactly: a proposal runs
without a card if and only if the function is granted to the chat, its
inputs validate, its scopes fit the grant, and $\ell_f \le \ell_c$.
At the default $\ell_c = 1$, an injected proposal can at most read and
make ordinary changes through granted agents; every wider change and
every action leaving the platform becomes a card. At $\ell_c = 0$,
only reads run unasked. At $\ell_c = 3$, every granted function runs
unasked---the case the trust level exists to make a person's explicit
choice. In every case, egress remains restricted to declared hosts and
no credential is exposed to the model.

\paragraph{Blast radius}
A compromised worker holds what its call was handed, its own home, and
its declared hosts. A compromised runtime can act for the people whose
delegations it holds---people with a chat open or a schedule
running---within 40 actions, and, since the level gate runs in the
runtime, could run their granted functions without opening a card; it
cannot answer a card, read the database, open a sealed value, obtain a
refresh token, or act for anyone else. Delegations last at most an hour
and are withdrawn by the person's stop-everything switch. A
compromised backend is a compromise of the deployment; it is the one
component trusted with everything.

\subsection{Documentation held to the code}
\label{sec:eval-docs}

This paper is distilled from the project's documentation, which names
the files implementing each mechanism. Five tests hold that
documentation to the code: every relative link and heading anchor
resolves; every backticked repository path exists; every setting the
configuration page describes is read by code, and every setting the
example files offer is documented; and the published list of actions is
regenerated from the backend's catalog and must match. These tests
check what a page \emph{points at}; whether a sentence is true is for a
reader to verify, which is why each mechanism here names its file.

\subsection{What is not evaluated}
\label{sec:eval-not}

\begin{itemize}
\item \textbf{Model behaviour.} No test calls a real model. The tests
  show that a proposal cannot pass the gates; they do not measure
  whether a given model resists injected text or chooses well. Neither
  is a claim of this paper.
\item \textbf{Performance and scale.} We do not report latency or
  throughput overheads, nor behaviour with many organizations, chats,
  or concurrent fires.
\item \textbf{The kernel.} Confinement results hold on the kernel
  tested; a kernel flaw is a way past every layer, and no test rules one
  out.
\item \textbf{The web app's own tests} exist but were not part of this
  evaluation.
\end{itemize}
